\section{Numerical Methods}

Numerical analysis studies mathematical and computational techniques designed
to obtain approximate solutions to mathematical problems. Such techniques are
particularly important when analytical solutions are unavailable, difficult to
obtain, or computationally impractical
\cite{burden2016numerical,chapra2021numerical}.

Classical numerical analysis literature identifies approximation, error,
convergence, stability, and computational efficiency as fundamental aspects in
the study and implementation of numerical algorithms
\cite{burden2016numerical,higham2002accuracy}.

One of the classical problems in numerical analysis is the determination of
roots of nonlinear equations. Given a function

\begin{equation}
    f(x) = 0,
\end{equation}

the objective is to determine a value $x^\ast$ that satisfies the equation, or
to compute an approximation sufficiently close to an exact root.

Several numerical algorithms may be employed for this purpose, each presenting
different convergence properties and computational costs. Among the methods
considered in this work are fixed-point iteration and the Newton--Raphson
method
\cite{burden2016numerical,chapra2021numerical}.

Fixed-point iteration begins by transforming the original problem

\begin{equation}
    f(x) = 0
\end{equation}

into an equivalent equation of the form

\begin{equation}
    x = g(x).
\end{equation}

Starting from an initial approximation $x_0$, a sequence is then constructed
according to

\begin{equation}
    x_{n+1} = g(x_n).
\end{equation}

If appropriate convergence conditions are satisfied, the sequence $\{x_n\}$
converges to a fixed point $x^\ast$ of the function $g$, satisfying

\begin{equation}
    g(x^\ast) = x^\ast.
\end{equation}

The convergence of fixed-point iteration depends on the properties of the
function $g$, the interval considered, and the initial approximation. A
frequently used sufficient condition involves the contraction properties of
$g$ in a neighborhood containing the fixed point
\cite{burden2016numerical}.

The recursive nature of this expression makes fixed-point iteration
particularly relevant to the present research. Its mathematical formulation
can be translated either into explicit iterative structures or into recursive
function definitions, thereby allowing a direct comparison between imperative
and functional programming approaches.

Another important method is the Newton--Raphson method. Starting from an
initial approximation $x_0$, successive approximations are computed using

\begin{equation}
    x_{n+1}
    =
    x_n
    -
    \frac{f(x_n)}{f'(x_n)}.
\end{equation}

Under appropriate assumptions, including sufficient smoothness of the function
and an initial approximation sufficiently close to a simple root, Newton's
method exhibits quadratic local convergence
\cite{burden2016numerical,chapra2021numerical}.

However, its practical behavior also depends on the function under analysis,
the selected initial approximation, and the numerical evaluation of both the
function and its derivative.

These methods illustrate a common characteristic of many numerical algorithms:
the repeated application of a mathematical procedure until a specified
criterion is satisfied. This property makes iterative numerical methods
particularly suitable for investigating how different programming models
represent repeated computation.
